\section*{Bloque B. Aplicaciones cuantitativas y optimización}

\noindent\textit{Preguntas 31 a 65. Requieren resolver economías de intercambio, calcular precios relativos, identificar la curva de contrato, evaluar mejoras de Pareto y aplicar funciones de bienestar social.}

\noindent\textit{Considere una economía de intercambio puro con dos consumidores A y B, dos bienes $x$ e $y$, dotaciones totales $X=10$ y $Y=10$. Las preferencias son $U_A=x_Ay_A$ y $U_B=x_By_B$. La dotación inicial es A$=(8,2)$ y B$=(2,8)$. Las preguntas 31 a 44 se refieren a esta economía.}

\q{La cantidad total del bien $x$ disponible en la caja de Edgeworth es:}
\begin{opts}
    \item 8
    \item 9
    \item 10
    \item 20
\end{opts}

\q{La cantidad total del bien $y$ disponible en la caja de Edgeworth es:}
\begin{opts}
    \item 8
    \item 9
    \item 12
    \item 10
\end{opts}

\q{La dotación inicial del consumidor A es:}
\begin{opts}
    \item $(2,8)$
    \item $(8,2)$
    \item $(5,5)$
    \item $(10,10)$
\end{opts}

\q{En la dotación inicial, la tasa marginal de sustitución del consumidor A es:}
\begin{opts}
    \item 4.00
    \item 2.00
    \item 1.00
    \item 0.25
\end{opts}

\q{En la dotación inicial, la tasa marginal de sustitución del consumidor B es:}
\begin{opts}
    \item 0.25
    \item 1
    \item 2
    \item 4
\end{opts}

\q{¿La dotación inicial es Pareto eficiente?}
\begin{opts}
    \item Sí, porque ambas tasas marginales de sustitución ya son iguales.
    \item No, porque difieren las RMS y hay ganancias de intercambio.
    \item Sí, porque las dotaciones totales de ambos bienes son fijas.
    \item No, porque los precios relativos iniciales son iguales a cero.
\end{opts}

\q{La curva de contrato de esta economía está dada por:}
\begin{opts}
    \item $y_A=2x_A$.
    \item $y_A=10-x_A$.
    \item $y_A=0$.
    \item $y_A=x_A$.
\end{opts}

\q{El precio relativo de equilibrio $p_x/p_y$ es:}
\begin{opts}
    \item 0.5
    \item 1
    \item 2
    \item 4
\end{opts}

\q{En el equilibrio competitivo, la cantidad del bien $x$ que consume A es:}
\begin{opts}
    \item 4
    \item 4.5
    \item 5
    \item 6
\end{opts}

\q{En el equilibrio competitivo, la cantidad del bien $y$ que consume A es:}
\begin{opts}
    \item 4
    \item 4.5
    \item 5
    \item 6
\end{opts}

\q{La utilidad de A en la dotación inicial es 16 y en el equilibrio es 25. La ganancia de utilidad de A es:}
\begin{opts}
    \item 0
    \item 5
    \item 7
    \item 9
\end{opts}

\q{La ganancia de utilidad del consumidor B al pasar de la dotación al equilibrio es:}
\begin{opts}
    \item 0
    \item 5
    \item 9
    \item 16
\end{opts}

\q{El equilibrio competitivo es Pareto eficiente porque:}
\begin{opts}
    \item A consume más que B.
    \item B pierde con el intercambio.
    \item se cumple $MRS_A=MRS_B=1$.
    \item los precios absolutos son iguales.
\end{opts}

\q{Si se desea alcanzar la asignación eficiente A$=(6,6)$ y B$=(4,4)$ con $p_x=p_y=1$, la transferencia lump-sum de B hacia A debe ser:}
\begin{opts}
    \item 2
    \item 1
    \item 4
    \item 0
\end{opts}

\q{Considere la misma economía, pero con dotación inicial A$=(9,1)$ y B$=(1,9)$. Si $p_x/p_y=1$, la cantidad de $x$ que consume A en equilibrio es:}
\begin{opts}
    \item 4
    \item 6
    \item 8
    \item 5
\end{opts}

\q{Con la dotación A$=(9,1)$ y B$=(1,9)$, si se transfiere 2 de B hacia A, la cantidad de $x$ que consume A es:}
\begin{opts}
    \item 4
    \item 5
    \item 6
    \item 7
\end{opts}

\q{La asignación A$=(6,6)$ y B$=(4,4)$ es Pareto eficiente porque:}
\begin{opts}
    \item las RMS de ambos consumidores son iguales.
    \item el consumidor B recibe menor cantidad que A.
    \item la dotación agregada permanece fija en el tiempo.
    \item los precios relativos de equilibrio valen cero.
\end{opts}

\q{El Segundo Teorema del Bienestar implica que:}
\begin{opts}
    \item toda asignación eficiente se sostiene con precios y transferencias lump-sum.
    \item todo equilibrio competitivo garantiza automáticamente equidad distributiva.
    \item la eficiencia asignativa exige necesariamente propiedad estatal de empresas.
    \item las transferencias lump-sum distorsionan los precios relativos de mercado.
\end{opts}

\q{En una caja de Edgeworth con $X=20$ e $Y=10$, si A tiene $(12,4)$, la cantidad de $x$ que tiene B es:}
\begin{opts}
    \item 4
    \item 6
    \item 10
    \item 8
\end{opts}

\q{En la misma caja, si A tiene $(12,4)$, la cantidad de $y$ que tiene B es:}
\begin{opts}
    \item 4
    \item 6
    \item 8
    \item 10
\end{opts}

\q{Si en una asignación interior con preferencias convexas se cumple $MRS_A=MRS_B=2$, entonces:}
\begin{opts}
    \item la asignación es Pareto eficiente.
    \item la asignación es ineficiente.
    \item solo es eficiente si los precios son iguales.
    \item solo es eficiente si las dotaciones son iguales.
\end{opts}

\q{Si $MRS_A=3$ y $MRS_B=1$, entonces:}
\begin{opts}
    \item no existen ganancias posibles derivadas del intercambio.
    \item la asignación actual ya es estrictamente Pareto eficiente.
    \item solo una expropiación directa mejora el bienestar conjunto.
    \item hay ganancias de intercambio al diferir las valoraciones.
\end{opts}

\q{Con dos bienes, $p_x=p_y=1$ y un exceso de demanda $z_x=4$, la ley de Walras implica que $z_y$ es:}
\begin{opts}
    \item $-4$
    \item $+0$
    \item $+4$
    \item $+8$
\end{opts}

\q{Si $p_x=p_y=1$ y $z_x=12$, entonces $z_y$ es:}
\begin{opts}
    \item $+12$
    \item $+00$
    \item $-12$
    \item $+24$
\end{opts}

\q{Si $p_x/p_y=2$ y $p_y=1$, entonces $p_x$ es:}
\begin{opts}
    \item 0.5
    \item 1
    \item 2
    \item 4
\end{opts}

\q{Si $p_x=3$ y $p_y=6$, el precio relativo $p_x/p_y$ es:}
\begin{opts}
    \item 2.0
    \item 0.5
    \item 3.0
    \item 6.0
\end{opts}

\q{Si la frontera de posibilidades de utilidad es $U_A+U_B=100$, el máximo valor de $U_A$ es:}
\begin{opts}
    \item 0
    \item 50
    \item 100
    \item 200
\end{opts}

\q{Si $W=\min(U_A,U_B)$ y $U_A+U_B=100$, el máximo de $W$ se alcanza con $U_A$ igual a:}
\begin{opts}
    \item 0
    \item 25
    \item 50
    \item 100
\end{opts}

\q{Si $W=U_A+U_B$ y $U_A+U_B=100$, el máximo de $W$ es:}
\begin{opts}
    \item 50
    \item 0
    \item 200
    \item 100
\end{opts}

\q{Una función de bienestar social que maximiza la utilidad del individuo peor situado corresponde a:}
\begin{opts}
    \item una función utilitarista.
    \item una función rawlsiana.
    \item el criterio de Pareto.
    \item el criterio de Kaldor-Hicks.
\end{opts}

\q{El Primer Teorema del Bienestar requiere:}
\begin{opts}
    \item mercados completos, sin externalidades y precio-aceptantes.
    \item presencia de monopolio natural y barreras legales a la entrada.
    \item existencia de externalidades negativas y bienes públicos puros.
    \item precios máximos uniformes fijados directamente por el Estado.
\end{opts}

\q{Si existe una externalidad negativa, el equilibrio competitivo:}
\begin{opts}
    \item es siempre Pareto eficiente.
    \item no es Pareto eficiente.
    \item es necesariamente equitativo.
    \item tiene exceso de demanda cero.
\end{opts}

\q{Un impuesto pigouviano óptimo debe igualar:}
\begin{opts}
    \item el costo externo marginal.
    \item el costo medio total.
    \item el costo total del bien.
    \item cero.
\end{opts}

\q{La condición de Samuelson para la provisión eficiente de un bien público es:}
\begin{opts}
    \item $MRS_A = MRS_B$
    \item $MRS_i = MRT_i$
    \item $\sum MRS = MRT$
    \item $\sum MRS_i = 0$
\end{opts}

\q{En un monopolio con $P>CMg$, el resultado es:}
\begin{opts}
    \item eficiencia asignativa plena en todos los mercados.
    \item ausencia absoluta de pérdida irrecuperable de bienestar.
    \item pérdida de bienestar y falla del Primer Teorema.
    \item un equilibrio competitivo Pareto eficiente sin costos.
\end{opts}
